%-------------------------
% Resume in Latex
% Author : Zorian Thornton
% Based off of: https://github.com/sb2nov/resume
% License : MIT
%------------------------

\documentclass[letterpaper,11pt]{article}

\usepackage{latexsym}
\usepackage[empty]{fullpage}
\usepackage{titlesec}
\usepackage{marvosym}
\usepackage[usenames,dvipsnames]{color}
\usepackage{verbatim}
\usepackage{enumitem}
\usepackage[hidelinks]{hyperref}
\usepackage{fancyhdr}
\usepackage[english]{babel}
\usepackage{tabularx}
\input{glyphtounicode}


%----------FONT OPTIONS----------
% sans-serif
% \usepackage[sfdefault]{FiraSans}
% \usepackage[sfdefault]{roboto}
% \usepackage[sfdefault]{noto-sans}
% \usepackage[default]{sourcesanspro}

% serif
% \usepackage{CormorantGaramond}
% \usepackage{charter}


\pagestyle{fancy}
\fancyhf{} % clear all header and footer fields
\fancyfoot{}
\renewcommand{\headrulewidth}{0pt}
\renewcommand{\footrulewidth}{0pt}

% Adjust margins
\addtolength{\oddsidemargin}{-0.5in}
\addtolength{\evensidemargin}{-0.5in}
\addtolength{\textwidth}{1in}
\addtolength{\topmargin}{-.5in}
\addtolength{\textheight}{1.0in}

\urlstyle{same}

\raggedbottom
\raggedright
\setlength{\tabcolsep}{0in}

% Sections formatting
\titleformat{\section}{
  \vspace{-4pt}\scshape\raggedright\large
}{}{0em}{}[\color{black}\titlerule \vspace{-5pt}]

% Ensure that generate pdf is machine readable/ATS parsable
\pdfgentounicode=1

%-------------------------
% Custom commands
\newcommand{\resumeItem}[1]{
  \item\small{
    {#1 \vspace{-2pt}}
  }
}

\newcommand{\resumeSubheading}[4]{
  \vspace{-2pt}\item
    \begin{tabular*}{0.97\textwidth}[t]{l@{\extracolsep{\fill}}r}
      \textbf{#1} & #2 \\
      \textit{\small#3} & \textit{\small #4} \\
    \end{tabular*}\vspace{-7pt}
}

\newcommand{\resumeSubSubheading}[2]{
    \item
    \begin{tabular*}{0.97\textwidth}{l@{\extracolsep{\fill}}r}
      \textit{\small#1} & \textit{\small #2} \\
    \end{tabular*}\vspace{-7pt}
}

\newcommand{\resumeProjectHeading}[2]{
    \item
    \begin{tabular*}{0.97\textwidth}{l@{\extracolsep{\fill}}r}
      \small#1 & #2 \\
    \end{tabular*}\vspace{-7pt}
}

\newcommand{\resumeSubItem}[1]{\resumeItem{#1}\vspace{-4pt}}

\renewcommand\labelitemii{$\vcenter{\hbox{\tiny$\bullet$}}$}

\newcommand{\resumeSubHeadingListStart}{\begin{itemize}[leftmargin=0.15in, label={}]}
\newcommand{\resumeSubHeadingListEnd}{\end{itemize}}
\newcommand{\resumeItemListStart}{\begin{itemize}}
\newcommand{\resumeItemListEnd}{\end{itemize}\vspace{-5pt}}

%-------------------------------------------
%%%%%%  RESUME STARTS HERE  %%%%%%%%%%%%%%%%%%%%%%%%%%%%


\begin{document}

%----------HEADING----------
\begin{center}
    \textbf{\Huge \scshape Zorian Thornton} \\ \vspace{1pt}
    \small (+1) 804-316-1769 $|$ \href{mailto:zoriant15@gmail.com}{\underline{zoriant15@gmail.com}} $|$
    \href{https://zorian15.github.io}{\underline{zorian15.github.io}} $|$ Seattle, WA
\end{center}

%-----------PROFESSIONAL SUMMARY-----------
\section{Professional Summary}
\small{Computational biologist and machine learning researcher with 6+ years of experience developing predictive models for biological systems. Expertise in building end-to-end deep learning pipelines for large-scale sequence analysis, implementing Bayesian statistical methods, and translating complex biological data into actionable insights. Proven track record of developing production-grade tools for viral evolution modeling and protein phenotype prediction.}

%-----------TECHNICAL SKILLS-----------
\section{Technical Skills}
 \begin{itemize}[leftmargin=0.15in, label={}]
    \small{\item{
     \textbf{Languages}{: Python, R, SQL, Java, MATLAB, bash} \\
     \textbf{Machine Learning}{: PyTorch, TensorFlow, JAX, NumPyro, scikit-learn, supervised/unsupervised learning, deep learning, Bayesian inference} \\
     \textbf{Statistical Methods}{: MCMC, variational inference, hierarchical modeling, dimensionality reduction, model benchmarking} \\
     \textbf{Bioinformatics \& Computational Biology}{: Phylogenetic inference, sequence analysis, deep mutational scanning, Augur, Auspice, Biopython, evofr} \\
     \textbf{Data Science \& Visualization}{: pandas, numpy, scipy, matplotlib, seaborn, statistical modeling} \\
     \textbf{Infrastructure \& Tools}{: Git, SLURM, CUDA, parallel computing, Unix/Linux, high-performance computing}
    }}
 \end{itemize}

%-----------EXPERIENCE-----------
\section{Research \& Work Experience}
  \resumeSubHeadingListStart

    \resumeSubheading
      {Fred Hutchinson Cancer Research Center}{Seattle, WA}
      {Doctoral Researcher, Computational Biology (Advisor: Frederick Matsen IV)}{July 2021 -- June 2026}
      \resumeItemListStart
        \resumeItem{Built end-to-end PyTorch pipelines for viral fitness prediction from deep mutational scanning data, processing datasets of tens of thousands of protein variants with genotype-phenotype mapping models}
        \resumeItem{Developed biophysically-inspired neural network architectures that predict variant effects from experimental measurements, enabling accurate phenotype predictions across multiple viral protein phenotypes}
        \resumeItem{Created open-source Java simulation framework for modeling viral evolution under immune selection, integrating phylogenetic inference with phenotypic fitness landscapes}
        \resumeItem{Collaborated with experimental biology teams to validate computational predictions, resulting in 2 peer-reviewed publications and tools adopted by multiple research groups}
      \resumeItemListEnd
      
    \resumeSubheading
      {Adaptive Biotechnologies}{Seattle, WA}
      {Machine Learning Computational Biology Intern}{June 2022 -- Sept. 2022}
      \resumeItemListStart
        \resumeItem{Implemented and benchmarked semi-supervised learning methods for T-cell receptor specificity classification using high-throughput sequencing datasets containing millions of sequences}
        \resumeItem{Designed scalable data preprocessing pipelines in Python to handle complex biological sequence data, optimizing memory usage and runtime for production deployment}
        \resumeItem{Conducted systematic downsampling experiments to characterize model performance across varying data conditions, identifying optimal training set sizes for different classification tasks}
        \resumeItem{Developed visualization dashboards to communicate technical findings to cross-functional teams including immunologists, clinicians, and product managers}
    \resumeItemListEnd

    \resumeSubheading
      {Virginia Tech Department of Statistics}{Blacksburg, VA}
      {Undergraduate Researcher (Advisor: Leah Johnson)}{Aug. 2017 -- May 2019}
      \resumeItemListStart
        \resumeItem{Developed predictive environmental models integrating temperature-dependent biological parameters with disease transmission dynamics to forecast Bluetongue virus spread across climate scenarios}
        \resumeItem{Built Bayesian hierarchical models using MCMC methods to quantify uncertainty in environmental-disease relationships, incorporating vector life history traits and host-pathogen interactions}
        % \resumeItem{Published results in Parasites \& Vectors demonstrating temperature-transmission suitability predictions for livestock disease management}
      \resumeItemListEnd

    \resumeSubheading
      {Nielsen}{Chicago, IL}
      {Data Science Intern}{June 2018 -- Aug. 2018}
      \resumeItemListStart
        \resumeItem{Implemented statistical framework to identify and correct errors in scanned receipt data, building automated pipeline for data quality assurance}
        \resumeItem{Developed anomaly detection algorithms to flag inconsistencies in consumer purchase data at scale}
      \resumeItemListEnd

    \resumeSubheading
      {University of Washington Department of Genome Sciences}{Seattle, WA}
      {Undergraduate Researcher (Advisor: Jim Bruce)}{June 2017 -- Aug. 2017}
      \resumeItemListStart
        \resumeItem{Built probabilistic models to quantify protein-protein interaction competition in cross-linking mass spectrometry data}
        \resumeItem{Developed 3D visualization tools for analyzing inter-surface regions of cross-linked proteins, contributing to published proteomics software (Journal of Proteome Research)}
      \resumeItemListEnd

  \resumeSubHeadingListEnd

%-----------EDUCATION-----------
\section{Education}
  \resumeSubHeadingListStart
    \resumeSubheading
      {University of Washington}{Seattle, WA}
      {Ph.D. in Genome Sciences}{Sep. 2019 -- March 2026}
    \resumeSubheading
      {Virginia Tech}{Blacksburg, VA}
      {B.Sc. Statistics \& B.Sc. Computational Modeling and Data Analytics}{Aug. 2015 -- May 2019}
      \resumeItemListStart
        \resumeItem{Minors in Mathematics and Computer Science}
      \resumeItemListEnd
  \resumeSubHeadingListEnd

%-----------PUBLICATIONS-----------
\section{Selected Publications}
 \begin{itemize}[leftmargin=0.15in, label={}]
    \small{\item{
    \textbf{Thornton, Z.}, Tran, T., Figgins, M., et al. (2026). Simulating the genetic and antigenic evolution of viruses under selection pressure from host immunity. \textit{Manuscript submitted to Virus Evolution.} \\
    Yu, T.*, \textbf{Thornton, Z.*}, Hannon, W., et al. (2022). A biophysical model of viral escape from polyclonal antibodies. \textit{Virus Evolution}, 8(2), veac110. \\
    El Moustaid, F., \textbf{Thornton, Z.}, et al. (2021). Predicting temperature-dependent transmission suitability of bluetongue virus in livestock. \textit{Parasites \& Vectors}, 14(1), pp.1-14. \\
    Keller, A., Chavez, J.D., Eng, J.K., \textbf{Thornton, Z.} and Bruce, J.E. (2018). Tools for 3D Interactome Visualization. \textit{Journal of Proteome Research}, 18(2), pp.753-758.
    }}
 \end{itemize}

%-----------PRESENTATIONS-----------
\section{Selected Presentations}
 \begin{itemize}[leftmargin=0.15in, label={}]
    \small{\item{
    \textbf{Thornton, Z.} Predicting Disease Progression of Colorectal Cancer via Self-Organizing Maps. RECOMB 2019, George Washington University, Washington DC, May 2019. \\
    \textbf{Thornton, Z.} Modeling Bluetongue Virus via Markov Chain Monte Carlo Methods. Student Experiential Learning Conference, Virginia Tech, Blacksburg, VA, Apr. 2018. \\
    \textbf{Thornton, Z.} Viewing Molecular Interaction Interfaces Through Directed Computational Methods. Department of Genome Sciences Research Symposium, University of Washington, Seattle, WA, Aug. 2017.
    }}
 \end{itemize}

%-----------HONORS & AWARDS-----------
\section{Honors \& Awards}
 \begin{itemize}[leftmargin=0.15in, label={}]
    \small{\item{
    \textbf{National Science Foundation Graduate Research Fellowship Program}, Honorable Mention \hfill 2020 \\
    \textbf{Mu Sigma Rho}, National Honor Society for Statistics \hfill 2019 \\
    \textbf{College of Science Dean's Roundtable Scholarship}, Virginia Tech \hfill 2018--2019 \\
    \textbf{Luther and Alice Hamlett Undergraduate Research Grant}, Virginia Tech \hfill 2018--2019 \\
    \textbf{Fralin Undergraduate Research Fellowship}, Virginia Tech \hfill 2017--2018 \\
    \textbf{Eckert Statistics Scholar}, Virginia Tech \hfill 2016--2019
    }}
 \end{itemize}

%-------------------------------------------
\end{document}
